% Comparison synthetic phase-2 vs ShezhenV3 full validation
\begin{table}[t]
\centering
\caption{Evolution from synthetic phase-2 to ShezhenV3 full validation.}
\label{tab:synthetic-vs-tcm}
\small
\begin{tabular}{lcc}
\hline
Quantity & Synthetic (phase 2) & ShezhenV3 (full) \\
\hline
Calibration $n$ & 240 & 1{,}144 \\
Clear / blur medians & 0.819 / 0.191 & 0.564 / 0.431 \\
$\tau$ & 0.505 & 0.498 \\
M5 Spearman $\rho$ & $\approx$0.75 & 0.47 \\
Test frames / seed & 500 (sim.\ $Q$) & 500 (real scores) \\
B2 M1 p50 (ms) & 364.7 & 341.2 \\
B2 M2 valid rate & 0.988 & 1.000 \\
\hline
\end{tabular}
\end{table}

\paragraph{Reading the migration.}
Absolute $Q_{\mathrm{img}}$ drops on real texture, and Spearman correlation decreases because overlap between clear and blurred cohorts increases (see \S\ref{sec:tcm-iqa}).
Nevertheless, B2 retains and slightly strengthens the valid-frame advantage while improving median RTT versus B0, demonstrating that the routing architecture generalizes beyond synthetic histograms.
Future work will explore ROI-limited scoring inside ShezhenV3 bounding boxes to tighten cohort separation without abandoning interpretable features.

\paragraph{JSONL audit example.}
Supplementary Listing~1 shows a de-identified JSONL row from the Franka auxiliary track; fields match Table~IV in \S\ref{sec:closed-loop}:
\begin{verbatim}
{"trial_id":1,"image_path":"experiments/synthetic/...",
 "q_img":0.8075,"flags":[],"retry_count":0,
 "route_decision":"upload_cloud","latency_ms":15.89}
\end{verbatim}
Timestamps \texttt{t\_capture}, \texttt{t\_iqa}, and \texttt{t\_route} are omitted in the abbreviated listing but required in full logs validated by \texttt{scripts/validate\_jsonl.py}.

\paragraph{Compute budget.}
Full validation scoring (1{,}144 pairs) plus three test seeds ($3\times500$ frames $\times$4 baselines simulated) completes in under ten minutes on a laptop CPU without GPU, supporting frequent re-calibration when new TCM batches arrive.
